% BHU PYQ -- Transcribed & Corrected
% Paper: ESMJ21 / ESMN21: Mineral and Earth Material
% Exam: B.Sc. (Hons.) Semester II Examination, 2025-26 (NEP)
% Department: Department of Geology
% Paper Ref: Conf/Even/N/2025-26/51

\documentclass[11pt,a4paper]{article}
\usepackage[a4paper,top=2.5cm,bottom=2.5cm,left=2.5cm,right=2.5cm,headheight=14pt]{geometry}
\usepackage{amsmath,amssymb}
\usepackage[T1]{fontenc}
\usepackage[utf8]{inputenc}
\usepackage{lmodern,microtype,fancyhdr,enumitem,hyperref,lastpage,parskip}

\hypersetup{colorlinks=true,urlcolor=black,linkcolor=black,pdftitle={ESMJ21: Mineral & Earth Material -- BHU Sem II 2025-26 (NEP)}}
\pagestyle{fancy}\fancyhf{}
\fancyhead[L]{\small   \textit{Banaras Hindu University}}
\fancyhead[R]{\small   \textit{ESMJ21 / ESMN21: Mineral \& Earth Material (2025-26)}}
\fancyfoot[C]{\small Page  \thepage\ of \pageref{LastPage}}


\newlist{parts}{enumerate}{2}
\setlist[parts,1]{label=(\alph*),leftmargin=*,itemsep=4pt,topsep=2pt}


\begin{document}


\begin{center}
  {\large\bfseries BANARAS HINDU UNIVERSITY}\\[2pt]
  {
\normalsize B.Sc. (Hons.) Semester II Examination 2025-26 (NEP)}\\[6pt]
  \rule{0.8\linewidth}{0.5pt}\\[6pt]
  {\large\bfseries DEPARTMENT OF GEOLOGY}\\[4pt]
  {\large\bfseries Paper Code: ESMJ21 / ESMN21 --- Mineral and Earth Material}\\[4pt]
  {
\normalsize Time: 2:30 Hours\hfill Maximum Marks: 70}\\[2pt]
  {\small Ref Code: Conf/Even/N/2025-26/51}
\end{center}

\rule{\linewidth}{0.4pt}\smallskip



\noindent   \textit{Instruction: Write your Roll No. at the top immediately on the receipt of this question paper.}\\[2pt]



\noindent   \textit{Note: Answer five questions in all, selecting two questions each from Section--A and Section--B. Section--C is compulsory. Figures in the right-hand margin indicate marks.}
\bigskip

\section*{SECTION -- A}




\noindent   \textbf{Question 1} \hfill [14 Marks]\\
With the help of a neat and well-labelled diagram, describe the components of a polarizing (petrological) microscope and explain the function of each component.

\bigskip




\noindent   \textbf{Question 2} \hfill [14 Marks]\\
Answer any   \textbf{two} of the following ($2  \times 7 = 14$ Marks):

\begin{parts}
  \item Briefly describe the method of rock thin-section preparation.
  \item Discuss the classification of silicate group of minerals.
  \item Differentiate between the physical properties of feldspar and calcite.
\end{parts}

\bigskip




\noindent   \textbf{Question 3} \hfill [14 Marks]\\
Describe the different mineral-forming processes. Discuss each process with suitable examples of minerals formed under them.

\bigskip




\noindent   \textbf{Question 4} \hfill [14 Marks]\\
Answer the following questions with suitable examples ($4  \times 3.5 = 14$ Marks):

\begin{parts}
  \item Why do diamond and graphite have the same chemical composition but different specific gravities?
  \item Explain why a mineral with cleavage can also show fracture.
  \item Why is diamond not just ten times harder than talc on the Mohs hardness scale?
  \item Explain how colour can be a diagnostic property in the identification of certain minerals.
\end{parts}

\bigskip
\section*{SECTION -- B}




\noindent   \textbf{Question 5} \hfill [14 Marks]\\
Hornblende and Orthopyroxene both possess prismatic crystal faces. Do they belong to the same crystallographic system? Justify your answer. Differentiate between these two minerals based on their optical properties. Is there any difference between a prismatic face and a pyramidal face?

\bigskip




\noindent   \textbf{Question 6} \hfill [14 Marks]\\
Explain in brief with suitable examples ($2  \times 7 = 14$ Marks):

\begin{parts}
  \item Is it always necessary that parallel crystal faces would have similar Miller indices?
  \item Is it possible to deduce Weiss parameters from Miller indices?
\end{parts}

\bigskip




\noindent   \textbf{Question 7} \hfill [14 Marks]\\
How would you distinguish between optically isotropic and optically anisotropic minerals? Provide at least one example each of optically isotropic and anisotropic minerals and also mention their crystal system. Two Olivine crystals (a common gemstone: 'Peridot') excavated from the basaltic rocks of two different depths from an igneous complex -- do they possess the same interfacial angles? Justify your answer with suitable reasons.

\bigskip
\section*{SECTION -- C (Compulsory)}




\noindent   \textbf{Question 8} \hfill [14 Marks]\\
Define / Answer the following in 1 word or sentence each:

\begin{parts}
  \item Name a mineral which has different colour and streak. \hfill [1 Mark]
  \item Name a tectosilicate mineral which contains cleavage. \hfill [1 Mark]
  \item Name a mineral that commonly exhibits third-order interference colours in a standard thin section. \hfill [1 Mark]
  \item Which mineral shows bladed habit? \hfill [1 Mark]
  \item Name a mineral that commonly shows polysynthetic twinning. \hfill [1 Mark]
  \item Define the following ($9  \times 1 = 9$ Marks):
  
\begin{enumerate}[label=(\roman*)]
    \item Mineral
    \item Plane of symmetry
    \item Axis of rotation
    \item Solid angle
    \item Euler's formula
    \item Birefringence
    \item Extraordinary ray
    \item Nicol prism
    \item Silicate polymerization
  \end{enumerate}
\end{parts}

\end{document}
